\section{Introduction}
\label{sec:introduction}

% -------------------------------------------------------------------------
\subsection{Background}
% -------------------------------------------------------------------------

Vehicle detection is a foundational task in intelligent transportation systems (ITS).
Accurate, real-time identification of vehicle types enables a wide range of applications,
including traffic-flow estimation, parking management, road-capacity planning, and
automated enforcement \cite{TODO-ITS-survey}.
Camera-based detection has become the preferred approach because a single fixed or
mobile camera can substitute for multiple inductive-loop sensors while simultaneously
capturing richer information such as vehicle class, speed, and trajectory
\cite{TODO-camera-vs-loop}.

Deep convolutional neural networks have transformed object detection since the
introduction of region-based convolutional networks \cite{TODO-rcnn}.
Among detector families, the \textit{You Only Look Once} (YOLO) series has gained
particular traction in real-time applications because its single-stage design
processes the entire image in one forward pass, achieving a favourable trade-off
between accuracy and inference speed \cite{TODO-yolo-survey}.
Successive YOLO generations have improved backbone efficiency, neck design, and
training strategies while keeping inference costs manageable on commodity hardware.

A recent architectural trend abandons non-maximum suppression (NMS) during inference.
NMS is a post-processing step that removes duplicate bounding boxes, but it introduces
additional latency and requires a manually tuned score threshold that affects
precision-recall balance \cite{TODO-nms-free}.
NMS-free detectors instead assign predictions during training so that only one
prediction per object survives at inference time.
The YOLO26 architecture, used in this study, adopts a dual-head assignment strategy
to achieve NMS-free inference together with the removal of Distribution Focal Loss and
the introduction of the MuSGD optimiser \cite{TODO-yolo26}.
We use the medium-scale variant YOLO26m (21.78 GFLOPs, 74.70\,M parameters, 42.0\,MB
weights), as confirmed by the model-size measurements in our experiment harness.

% -------------------------------------------------------------------------
\subsection{Problem Statement and Motivation}
% -------------------------------------------------------------------------

Despite strong benchmark performance on large datasets such as COCO, two practical
problems limit how well published results transfer to real traffic deployments.

\textbf{Domain shift.}
Traffic composition and scene geometry vary substantially across regions.
South and Southeast Asian traffic contains vehicle types that are rare or absent in
Western benchmarks, including buses of varying sizes, vans used as informal taxis, and
motorcycles mixed with heavy goods vehicles \cite{TODO-domain-shift}.
Detectors trained on generic datasets typically lose accuracy when applied to such
scenes, yet many applied studies train and test on a single domain and do not measure
this loss.
The dataset used here covers three classes — \textit{bus}, \textit{car}, and
\textit{van} — assembled from a heterogeneous collection of scenes and exported from
Roboflow Universe under CC~BY~4.0 (project
\texttt{vehicle-detection-vznzd-dkl8g},~v1, 9\,725 images in total)
\cite{roboflow_dataset}.
Class boundaries, particularly between \textit{car} and \textit{van}, are
visually ambiguous, as confirmed by our TIDE-style error analysis which found
\textit{van}$\rightarrow$\textit{car} to be the dominant confusion pair.

\textbf{Silent label corruption in mixed-format datasets.}
Community datasets assembled on annotation platforms are often merged from several
sources and may contain annotations stored in different formats.
In Ultralytics~8.x, a label file that contains both bounding-box rows (five values per
row) and polygon rows (more than five values per row) is parsed entirely as a polygon
file; bounding-box rows are silently reinterpreted as degenerate polygons, producing
incorrect boxes without any error message \cite{TODO-ultralytics-label-bug}.
Our audit of the training split identified three label files in which this condition
holds (81 polygon rows mixed among 18\,413 bounding-box rows), all in the
\textit{bus} class.
The validation and test splits contained no mixed files.
The magnitude of the accuracy impact depends on how many ground-truth boxes are
corrupted; this study quantifies that effect.

% -------------------------------------------------------------------------
\subsection{Research Questions}
% -------------------------------------------------------------------------

This study addresses the following research questions.

\begin{enumerate}
\item[\textbf{RQ1.}]
  How accurate and how fast is YOLO26m for vehicle detection and counting on a
  heterogeneous public dataset, and how does it compare with YOLOv8m, YOLO11m, and
  RT-DETR-l under identical training settings?
  [DATA DIBUTUHKAN: Tables T8, T9 — baseline and augmentation ablation runs not yet
  executed.]

\item[\textbf{RQ2.}]
  How common are mixed bounding-box/polygon label files in public vehicle datasets,
  and how much does correcting them change detection accuracy?

\item[\textbf{RQ3.}]
  How suitable is the trained detector for deployment in a low-cost Streamlit
  pipeline on consumer hardware (Apple M4 Pro), and what end-to-end latency does
  it achieve?
\end{enumerate}

% -------------------------------------------------------------------------
\subsection{Contributions}
% -------------------------------------------------------------------------

The specific contributions of this paper are as follows.

\begin{enumerate}
\item \textbf{Benchmarking YOLO26 on a heterogeneous traffic dataset.}
  We evaluate YOLO26m under a reproducible training protocol and compare it against
  three competitive baselines (YOLOv8m, YOLO11m, RT-DETR-l).
  On the test split, the deployed YOLO26m checkpoint achieves
  P\,=\,0.845, R\,=\,0.752, mAP$_{50}$\,=\,0.784, and
  mAP$_{50:95}$\,=\,0.594 (clean labels, as measured by our experiment harness).
  [DATA DIBUTUHKAN: comparative tables once baselines are trained.]

\item \textbf{Documentation and quantification of a silent label-format bug.}
  We identify, reproduce, and measure the accuracy effect of mixed bounding-box/polygon
  label files in Ultralytics exports, and release an open-source audit script
  (\texttt{trace\_experiments.py}, stage \texttt{audit}) that other practitioners can
  apply to their own datasets.

\item \textbf{TIDE-style error taxonomy.}
  Using our TIDE-inspired analysis (confidence threshold 0.4, IoU$_{fg}$\,=\,0.5,
  IoU$_{bg}$\,=\,0.1), we decompose test-set errors into classification, localisation,
  duplicate, background, and missed categories per class and per object size,
  providing actionable insight for future model improvements.

\item \textbf{End-to-end deployment pipeline and latency measurement.}
  We package the trained detector in a Streamlit application that logs vehicle counts
  to Google Sheets via an n8n webhook and sends Telegram alerts.
  Measured CPU inference latency is 96.3\,ms per image (mean end-to-end: 99.1\,ms);
  MPS (Apple M4 Pro GPU) latency is [DATA DIBUTUHKAN: MPS run not yet captured in
  summary].
\end{enumerate}

% -------------------------------------------------------------------------
\subsection{Paper Structure}
% -------------------------------------------------------------------------

The remainder of this paper is organised as follows.
Section~\ref{sec:related_work} reviews related work on vehicle detection, NMS-free
detectors, and label-quality issues in crowdsourced datasets.
Section~\ref{sec:methodology} describes the dataset, the label audit procedure, the
model architecture, and the deployment pipeline.
Section~\ref{sec:experiments} defines the experimental protocol, evaluation metrics,
and hardware configuration.
Section~\ref{sec:results} presents the detection performance, latency measurements,
and error analysis.
Section~\ref{sec:discussion} interprets the results with respect to the research
questions and discusses limitations.
Section~\ref{sec:conclusion} summarises the findings and outlines future work.
